\section{Locks}

\subsection{Hardware Support for Parallelism}

\begin{remark}
    \textbf{Software algorithms} such as Bakery or Filter locks have \textbf{linear space complexity} $O(n)$, making them impractical for large numbers of processes. Hardware support is needed to implement locks efficiently.
\end{remark}

\begin{remark}
    Moder architectures offer \textbf{atomic instructions} that allow for operations to happen in a single atomic step. This allows mutex-implementation with \textbf{constant space complexity} $O(1)$ and can be used for lock-free programming.
\end{remark}

\subsubsection{Typical Hardware Instructions}

\begin{definition}
    \textbf{Compare and Swap (CAS):} Atomically compares the value of a memory location with a given value. If they are equal, it updates the memory location with a new value. Otherwise, it does nothing.
\end{definition}

\begin{definition}
    \textbf{Test and Set (TAS):} Atomically reads the value of a memory location and sets it to 1.
\end{definition}

\begin{definition}
    \textbf{Load Linked / Store Conditional (LL/SC): } LL reads the value of a memory location and sets a "link" to it. SC writes a new value to the memory location only if the link is still valid (i.e., no other process has written to that location since the LL).
\end{definition}

\begin{remark}
    \textbf{Java} offers these funcitonalities via the classes \textbf{AtomicInteger}, \textbf{AtomicLong}, etc.
\end{remark}

\subsubsection{Locks using atomic operations}

\begin{codeexample}
//TAS LOCK IN Java
public class TASLock implements Lock {
    AtomicBoolean state = new AtomicBoolean(false);

    public void lock() {
        while (state.getAndSet(true)) {
            //busy wait
        }
    }

    public void unlock() {
        state.set(false);
    }
}
\end{codeexample}

\begin{remark}
    This lock consistently checks the state variable with an atomic write operation, which causes a lot of cache coherence traffic. To avoid this we can use TATAS lock (checking if the state is false before trying to acquire the lock).
\end{remark}

\begin{codeexample}
//TAS LOCK IN Java
public class TATASLock implements Lock {
    AtomicBoolean state = new AtomicBoolean(false);

    public void lock() {
        do {
            while (state.get());
        } while (state.getAndSet(true));
    }

    public void unlock() {
        state.set(false);
    }
}
\end{codeexample}

\begin{remark}
    To further optimize this lock we can use a backoff strategy (e.g. exponential backoff) to reduce the cache coherence traffic. This means that if a process fails to acquire the lock, it will wait for a certain amount of time (ex. exponentially increasing) before trying again.
\end{remark}

\subsection{Locking algorithms}

To implement a Lock (Mutex) correctly using only basic memory reads/writes, three conditions must be met:
\begin{enumerate}
    \item \textbf{Mutual Exclusion:} No two processes can be in the critical section (CS) simultaneously.
    \item \textbf{Freedom from Deadlock:} If processes try to enter the CS, at least one must eventually succeed.
    \item \textbf{Freedom from Starvation:} Every process that attempts to enter must eventually succeed.
\end{enumerate}

\subsubsection{State Space Diagrams}

\begin{definition}
    A \textbf{State Space Diagram} is a formal tool used to track every possible configuration of a concurrent program.
    \textbf{States} are typically represented as a tuple, e.g., $[p, q, \text{wantp}, \text{wantq}]$, where $p$ and $q$ are the current program counters (line numbers) of the threads.
\end{definition}

\begin{remark}
    We can use state space diagrams to prove the correctness of lock implementations and understand the follwoing algorithms and idenfy failures
    \begin{itemize}
        \item \textbf{No Mutex:} Reachable state where both processes are in the CS.
        \item \textbf{Deadlock:} A state with no outgoing transitions where threads are stuck in the protocol.
        \item \textbf{Starvation:} An execution path where one thread never reaches the CS despite trying.
    \end{itemize}
\end{remark}

\subsubsection{Attempts to mutual exclusion}

\begin{algorithm}
volatile int flag[2] = {0, 0};

// Thread 0
while (flag[1] != 1) {}
flag[0] = 1;
// Critical section
flag[0] = 0;
\end{algorithm}

Observe that this algorithm fails, as both thread might check the others flag and see that it is 0, and then both enter the critical section.

\begin{algorithm}
volatile int flag[2] = {0, 0};

// Thread 0
flag[0] = 1;
while (flag[1] != 1) {}
// Critical section
flag[0] = 0;
\end{algorithm}

Observe that this algorithm also fails, as both thread might set their flag to 1 and then check the others flag and see that it is 1 and therefore deadlock.

\begin{algorithm}
volatile int turn = 1;

// Thread 0
while (turn == 1) {}
// Critical section
turn = 0;
\end{algorithm}

Observe that this algorithm also fails as we do not guarantee that a thread resets the turn variable if it crahses or does not use the critical section. Furthermore this algorithm would ensure strict alternation (round-robin) between the threads.


\subsubsection{Successful Algorithms}

The first correct solution to the mutual exclusion problem for two processes. It combines the flag system with a \textbf{turn} variable to break ties. If both want to enter, the \textbf{turn} variable decides who must temporarily retract their interest.

\begin{codeexample}
// Process P
wantp = true;
while (wantq) {
    if (turn == 2) {
        wantp = false;
        while (turn == 2); // wait for turn
        wantp = true;
    }
}
// Critical Section
turn = 2;
wantp = false;
\end{codeexample}

Peterson's algorithm is a more concise version of Dekker's. Instead of retracting interest, a process simply "gives way" by setting a \textbf{victim} variable to itself.

\begin{codeexample}
// Process P
wantp = true;
victim = 1; // P is process 1
while (wantq && victim == 1);
// Critical Section
wantp = false;
\end{codeexample}

\begin{theorem}
    Peterson's Algorithm provides Mutual Exclusion, Deadlock Freedom, and Starvation Freedom, and works elegantly using only atomic registers.
\end{theorem}

\begin{definition}
    An event $A$ is said to precede $B$ denoted as $A \prec B$ if the execution interval of $A$ does not overlap with the execution interval of $B$. If they do not overlap, then $A$ and $B$ are called concurrent.
\end{definition}

\begin{definition}
    An \textbf{atomic register} is a register that performs operations at a single point in time hence two operations on the same register always have diffrent effect times.
\end{definition}

\begin{proof}
    \textbf{Mutual Exclusion of Peterson's Algorithm:} By contradiction: assume concurrent $CS_P$ and $CS_Q$. Assume wlog: 
    $$W_Q(\text{victim}=Q) \rightarrow W_P(\text{victim}=P)$$

    From the program code:
    \begin{itemize}
        \item $W_P(\text{flag}[P]=\text{true}) \rightarrow W_P(\text{victim}=P) \rightarrow R_P(\text{flag}[Q]) \rightarrow R_P(\text{victim}) \rightarrow CS_P$
        \item $W_Q(\text{flag}[Q]=\text{true}) \rightarrow W_Q(\text{victim}=Q) \rightarrow R_Q(\text{flag}[P]) \rightarrow R_Q(\text{victim}) \rightarrow CS_Q$
    \end{itemize}

    Since $W_Q(\text{victim}=Q) \rightarrow W_P(\text{victim}=P)$, thread $P$ must read $P$ for victim.
    For thread $P$ to enter its critical section ($CS_P$), the condition while condition must be false.

    Because $victim == P$ is true, it follows that $R_P(\text{flag}[Q])$ must have read \textbf{false}.
    However, we have the sequence: 
    $$W_Q(\text{flag}[Q]=\text{true}) \rightarrow W_Q(\text{victim}=Q) \rightarrow W_P(\text{victim}=P) \rightarrow R_P(\text{flag}[Q])$$
    By the transitivity of the precedence relation $\rightarrow$, $R_P(\text{flag}[Q])$ must read \textbf{true}.
    This is a contradiction. Therefore, mutual exclusion is satisfied.
\end{proof}

\begin{proof}
\textbf{Proof: Freedom from Starvation:} Assume wlogthat thread $P$ runs forever in its lock loop, waiting until flag[Q] is flase or victim is not P.

There are three possibilities for thread $Q$:
\begin{enumerate}
    \item \textbf{$Q$ is stuck in its non-critical section:} 
    Then flag[Q] = false, and $P$ can continue. $\implies$ Contradiction.
    
    \item \textbf{$Q$ repeatedly enters and leaves its critical section:} 
    Each time $Q$ enters, it sets $victim = Q$. Since $P$ does not change victim again, victim remains $Q$, allowing $P$ to continue. $\implies$ Contradiction.
    
    \item \textbf{$Q$ is also stuck in its lock loop:} 
    $Q$ would be waiting until flag[P] is false or victim is not Q. However, the conditions victim == P and victim == Q cannot hold at the same time. $\implies$ Contradiction.
\end{enumerate}
In all cases, we reach a contradiction, proving the algorithm is starvation-free.
\end{proof}

\begin{remark}
    The Peterson's algorithm is correct, but it is not scalable. It is only proven for 2 processes. We can use the so called \textbf{Filter Lock} where every process has a level and in order to enter the critical section a process has to go trough all lower levels where on each level we use peterson's algorithm to eliminate one victim and let the others pass. This ensures mutual exclusion.
\end{remark}

\begin{remark}
    Another algorithm for mutual exclusion of $n$ processes is the \textbf{Bakery Algorithm}. Every process has to take a numbered ticket with a number greater than all outstanding tickets and waits until it has the lowest number.
\end{remark}

\begin{theorem}
    If $S$ is a [atomic] read/write system with at least two processes and $S$ solves mutual exclusion with global progress [Deadlock freedom], then $S$ must have at leas as many variables as processes.
\end{theorem}

\begin{remark}
    But how can there exist computers with thousands to millions of threads? (Solution: Mutex is not implemented as we thought)
\end{remark}


\subsection{Locking Strategies} 

\subsubsection{Fine-Grained Locking}

\begin{remark}
    Instead of locking everything coarse-grained (the whole data structure) we can lock individual parts of it. This is called \textbf{fine-grained locking} and allows for higher concurrency.
\end{remark}

\subsubsection{Read-Write Locks}

\begin{remark}
    With coarse- and fine-grained locking we have complete mutual exclusion on resources but what if we only need exclusive access for writing?
\end{remark}

\begin{definition}
    A \textbf{read-write lock} allow fair access to resources by allowing one writer (every N reads) or multiple readers at the same time.
\end{definition}

\begin{remark}
    To implement a RWL we need to keep track of the number of readers waiting, the number of writers waiting and the number of readers a writer has to wait until it can write (this ensures fairness).
\end{remark}

\subsubsection{Optimistic Locking}

\begin{remark}
    When for example traversing a list, finding a node with \textbf{hand-over-hand locking} requires locking all nodes on the way which consumes time and energy. But how high is it that two threads edit the same node at the same time? 
\end{remark}

\begin{definition}
    \textbf{Optimistic Locking: } Instead of locking every node we could traverse the list without locking and only lock the node we are interested in. If we lock and notice that the node has been modified and our locks on (predecessor/successor) are nolonger valid we restart the operation.
\end{definition}

\begin{remark}
    Optimistic locking allows us for less overhead on locking but might require multiple traversals if there is high contention.
\end{remark}

\subsubsection{Lazy Locking}

\begin{definition}
    \textbf{Lazy locking} is a technique where data-elements are not immediately updated, but rather marked as deleted (or inserted) and only updated when a thread traverses the data structure again. This saves the overhead of multiple traversals.
\end{definition}

\subsubsection{Lock-Free Data Structures}

\begin{remark}
    Using atomic operations (mostly CAS) we can implement data structures without locks. 
\end{remark}

\begin{codeexample}
class ConcurrentCounter {
    private AtomicLong value;
    
    public ConcurrentCounter() {
        value = new AtomicLong(0);
    }
    
    public void increment() {
        long expected;
        do {
            expected = value.get();
        } while (!value.compareAndSet(expected, expected + 1));
    }
    
    public long getValue() {
        return value.get();
    }
}
\end{codeexample}

\begin{remark}
    A lock free algorithm is always deadlock free
\end{remark}

\begin{definition}
    \textbf{Lock-Free: } System wide progress is guaranteed. This means that at least one process can always make progress.
\end{definition}

\begin{definition}
    \textbf{Wait-Free: } Any operation on the data structure will terminate in a bounded number of steps.
\end{definition}

\subsubsection{The ABA Problem}

\begin{remark}
    Assume we have a stack that we try to implement lock free and to save computation we dont fully garbage collect (reclaim memory) of popped elements. Instead we store them in a pool.
\end{remark}

\begin{example}
   Assume we pop element $a$ from the stack and store it in the pool. Then we push some element $b$ and then $a$ again (which we got from the pool, hence has the same adress). A thread that read the last element before poping $a$ and now checks again will not notice the change of the stack.
\end{example}

\begin{important}
    The \textbf{ABA problem} is a common issue in lock-free programming where an atomic operation (like compare-and-swap) succeeds because the value has changed twice (e.g., A -> B -> A), making it appear unchanged to the operation. This can lead to incorrect program states.
\end{important}


\begin{definition}
    We can use \textbf{Tagged Pointers} to significantly reduce the probability of the ABA problem by incrementing a \textbf{tag} along with the pointer in the CAS operation to detect up to $2^k$ changes where $k$ is the number of bits used for the tag.
\end{definition}

\begin{definition}
    \textbf{Hazard Pointers:} A common technique to solve the ABA problem is to use hazard pointers. A hazard pointer is a pointer to a node that is currently being accessed by a thread. So when a thread reads the state of a datastructure it will store the address of the nodes it reads in its hazard pointer. Anyone trying to edit it will see the hazard pointer and not modify the nodes.
\end{definition}